% Figure: loss split and efficiency of the brushless DC machine against
% shaft torque at fixed speed. Copper loss rises as torque squared (since
% current is proportional to torque); core and friction losses are nearly
% fixed at fixed speed. Efficiency peaks where the fixed losses equal the
% copper loss, the same crossover seen in the transformer case study.
\begin{figure}[htbp]
\centering
\begin{tikzpicture}[>=Latex]
\begin{axis}[
  width=12cm, height=6.8cm,
  xlabel={shaft torque $\tau$ ($\si{\newton\meter}$)},
  ylabel={loss ($\si{\watt}$)},
  xmin=0, xmax=1.0,
  ymin=0, ymax=120,
  xtick={0,0.2,0.4,0.6,0.8,1.0},
  ytick={0,20,40,60,80,100,120},
  legend pos=north west,
  legend cell align=left,
  font=\scriptsize,
  grid=both,
  grid style={enggray!20},
  axis y line*=left,
]
% copper loss: I^2 R with I proportional to tau; at tau=0.9 Nm, I=12.5 A,
% R_phase summed gives copper loss ~ 95 W. Scale so P_cu = 117*tau^2.
\addplot[engred, line width=1pt, domain=0:0.95, samples=60] {117*x^2}
  \closedcycle;
\addlegendentry{copper loss $\propto \tau^2$}
% fixed loss (core + friction + windage) ~ 20 W at fixed speed
\addplot[engamber, line width=1pt, domain=0:0.95, samples=2] {20};
\addlegendentry{core + friction (fixed)}
% crossover where copper = fixed: tau = sqrt(20/117) = 0.41 Nm
\addplot[enggray, dashed] coordinates {(0.413,0) (0.413,40)};
\node[enggray, anchor=south, font=\scriptsize] at (axis cs:0.413,42)
  {$\tau \approx 0.41\,\si{\newton\meter}$};
\end{axis}
\begin{axis}[
  width=12cm, height=6.8cm,
  xmin=0, xmax=1.0,
  ymin=0, ymax=100,
  hide x axis,
  axis y line*=right,
  ylabel={efficiency $\eta$ (\%)},
  ytick={0,20,40,60,80,100},
  font=\scriptsize,
]
% efficiency: P_out = tau*omega with omega=314 rad/s; P_loss = 20 + 117 tau^2.
% eta = P_out/(P_out+P_loss). Plot in percent.
\addplot[engblue, line width=1.3pt, domain=0.02:0.95, samples=120]
  {100*(314*x)/(314*x + 20 + 117*x^2)};
\addlegendentry{efficiency (right axis)}
\legend{}
\node[engblue, anchor=north east, font=\scriptsize] at (axis cs:0.95,82)
  {efficiency};
\end{axis}
\end{tikzpicture}
\caption[Brushless DC machine loss split and efficiency]{Loss split and
efficiency of the case-study brushless DC machine against shaft torque at
the rated speed of $3000$ revolutions per minute. Copper loss rises as the
square of torque because phase current is proportional to torque; core
loss, friction, and windage are nearly fixed at fixed speed. Efficiency,
read on the right axis, climbs steeply out of no-load, where the fixed
losses dominate a small output, and peaks near the torque at which copper
loss equals the fixed loss, the same crossover that set the
peak-efficiency load of the transformer. Past the peak the rising copper
loss erodes efficiency again, which is why a machine that must run mostly
at light load is wound differently from one that runs near its rating.}
\label{fig:vol04ch08:bldc-efficiency}
\end{figure}
